%%%%%%%%%%%%%%%%%%%%%%%%%%%%%%%%%%%%%%%%%%%%%%%%%%%%%%%%%%%%%%%%%%%%%%%%%%%%
%% Draft manuscript for Operations Research / INFORMS style.
%% The template structure follows the INFORMS Operations Research author
%% submission template, Ver. 1.1, February 21, 2025.
%%%%%%%%%%%%%%%%%%%%%%%%%%%%%%%%%%%%%%%%%%%%%%%%%%%%%%%%%%%%%%%%%%%%%%%%%%%%

\documentclass[opre,dblanonrev]{informs4}
\usepackage{eqndefns-left}
\RequirePackage{tgtermes}
\RequirePackage{newtxtext}
\RequirePackage{newtxmath}
\RequirePackage{bm}
\RequirePackage{endnotes}

\OneAndAHalfSpacedXII

%%% OPRE uses endnotes.
\let\footnote=\endnote
\let\enotesize=\normalsize
\def\notesname{Endnotes}%
\def\makeenmark{\hbox to1.275em{\theenmark.\enskip\hss}}
\def\enoteformat{\rightskip0pt\leftskip0pt\parindent=1.275em
  \leavevmode\llap{\makeenmark}}

\usepackage{natbib}
\bibpunct[, ]{(}{)}{,}{a}{}{,}%
\def\bibfont{\small}%
\def\bibsep{\smallskipamount}%
\def\bibhang{24pt}%
\def\newblock{\ }%
\providecommand{\BIBand}{and}

\usepackage{booktabs}
\usepackage{enumitem}

\TheoremsNumberedThrough
\ECRepeatTheorems
\EquationsNumberedThrough

\newcommand{\Afull}{\mathcal A^{\mathrm{full}}}
\newcommand{\Acand}{\mathcal A^{\mathrm{cand}}}
\newcommand{\Ameas}{\mathcal A^{\mathrm{meas}}}
\newcommand{\Cfull}{\mathcal C^{\mathrm{full}}}
\newcommand{\Ccand}{\mathcal C^{\mathrm{cand}}}
\newcommand{\Real}{\mathbb R}
\newcommand{\E}{\mathbb E}
\newcommand{\one}{\mathbf 1}
\newcommand{\lowermu}{\underline{\mu}}
\newcommand{\normone}[1]{\left\|#1\right\|_1}

\MANUSCRIPTNO{OPRE-2026-XXXX}
\RUNAUTHOR{Anonymous}
\RUNTITLE{Robust Candidate MaxWeight Scheduling}

\TITLE{Robust Candidate MaxWeight Scheduling for Heterogeneous Compute Fabrics}

\ARTICLEAUTHORS{%
\AUTHOR{Anonymous Author(s)}
\AFF{Anonymous institution, \EMAIL{anonymous@example.com}}
}

\ABSTRACT{%
Heterogeneous compute clusters increasingly run workloads whose service rates
depend on co-location, host pressure, accelerator memory, and topology.  These
effects make the scheduler's feasible action space combinatorial, while the
service map is only partially measured.  We formulate this setting as a
configuration-action stochastic processing network and study a robust
candidate-restricted MaxWeight policy.  The theory separates three objects that
are often conflated in systems schedulers: the full configuration action space,
a computable candidate family, and a conservative lower-service map estimated
from profiling.  We prove that if the full-action support region has slack
\(\delta\), the candidate family has fabric-cover loss \(L\rho\), lower-service
estimation consumes \(\epsilon_{\mathrm{est}}\), queue-scaled penalties consume
\(\beta\), and the optimization oracle has queue-scaled error \(\alpha_1\), then
the remaining margin
\(\eta=\delta-(L\rho+\epsilon_{\mathrm{est}}+\beta+\alpha_1)>0\) yields a
finite-set Foster recurrence certificate under bounded conditional second
moments.  The proof is formalized in Lean for both fixed and statewise feasible
families.  We instantiate the framework in Scheduleurm, a production-oriented
research scheduler, using measured service slices and state-of-the-art
(SOTA)-style replay baselines on identical task lists.  On the current
finite-slice benchmark,
Scheduleurm improves all-job completion time over its legacy policy by
\(11.825\times\), \(1.079\times\), \(1.510\times\), and \(1.154\times\) on four
declared local or node-bucket workloads, and is not Pareto-dominated by the
SOTA-style baselines under the same measured service cache.  These empirical
claims are deliberately bounded: we compare policy semantics on measured
service rows, not direct executions of external schedulers, and we do not claim
global fabric-cover calibration beyond the audited finite slices.
}%

\KEYWORDS{stochastic processing networks; MaxWeight scheduling; heterogeneous clusters; graphics-processing-unit co-location; robust optimization; formal verification; queueing networks}
\HISTORY{Draft of June 11, 2026.}

\begin{document}
\maketitle

\section{Introduction}
\label{sec:introduction}

Schedulers for modern research clusters operate in a regime that is poorly
described by a fixed resource vector.  A reinforcement learning (RL) job can
retain near-solo progress when several copies share a graphics processing unit
(GPU), whereas a dense GPU kernel can lose finite-batch delay performance when
the same GPU is aggressively multiplexed.  Central processing unit (CPU)-heavy
evaluation jobs stress host cores and memory rather than accelerator compute,
and short control-plane jobs expose queueing overhead and fragmentation.  These
cases are not exceptions to a simple model; they are the model.  The scheduler
chooses global configurations whose service rates depend on the set of jobs,
their placement, co-location profile, and local hardware state.

The operations-research difficulty is that the natural action is a global
configuration.  Enumerating all placements and co-location patterns is
combinatorial, but reducing the problem to a small list of hand-picked
``sweet spots'' can invalidate stability claims.  A policy may look efficient on
a finite batch while silently using an infeasible profile, or it may be
throughput-optimal in a model that is disconnected from the measured service
map.  Our starting point is therefore to keep the full action family in the
mathematical statement and to treat the implemented candidate family as an
approximation whose loss must be paid from explicit slack.

We make this separation precise for a heterogeneous compute fabric.  Jobs are
partitioned into finitely many classes \(i\in\mathcal I\).  At time \(t\), the
queue vector is \(Q(t)\), the scheduler chooses a global action
\(a_t\in\Afull\), and the random service vector is
\(S(a_t,Z_t,\omega_t)\), where \(Z_t\) denotes a hardware or workload regime.
The expected service map \(\mu(a)\) is unknown and is represented in the
implemented scheduler by conservative lower-service rows \(\lowermu(a)\).  The
candidate family \(\Acand\) is not assumed to be the full action family.  Instead
we require a finite-feature fabric cover and account for its support-function
loss.

The main theoretical contribution is a robust candidate MaxWeight theorem with
complete slack accounting.  If an arrival vector \(\lambda\) lies inside the
downward-closed full-action service region with support slack \(\delta\), and if
candidate approximation, lower-service estimation, queue-scaled penalties, and
approximate optimization consume \(L\rho\), \(\epsilon_{\mathrm{est}}\),
\(\beta\), and \(\alpha_1\), respectively, then the policy has negative
Lyapunov drift outside a finite set whenever
\[
  \delta > L\rho+\epsilon_{\mathrm{est}}+\beta+\alpha_1.
\]
Under a concrete bounded conditional second-moment stochastic model, this drift
implies a finite-set Foster recurrence certificate.  The theorem is stated for
fixed feasible families and for statewise feasible families, the latter being
the case relevant to live schedulers whose feasible actions depend on available
jobs, alive nodes, and capacity-boundary certificates.

The empirical contribution is not a claim that a single implementation has
solved cluster scheduling in full generality.  We instantiate the theorem in
Scheduleurm, an existing scheduler used for local research workloads, and report
three bounded pieces of evidence.  First, measured service curves close declared
finite slices for four workload buckets: local light-control tasks, GPU-heavy
JAX numerical-computing matrix-multiplication tasks, local CPU-heavy tasks, and
Robust-Ensemble Soft Actor-Critic (RE-SAC) / Bayesian Amnesic Piecewise-Robust
(BAPR)-like hybrid RL tasks.  Second, explicit task-list replay compares the
current policy with the legacy policy and with SOTA-style policy semantics, all
using the same measured service cache.  Third, production and live-trace audits
show that the lower-service oracle bridge is executable on the current
completed-active Scheduleurm population and on one emitted live scheduler trace.
These results are useful because they connect the proof obligations to scheduler
artifacts; they do not replace broader perturbation profiling for a global
fabric-cover constant.

The paper is organized as follows.  Section~\ref{sec:model} introduces the
configuration-action processing network.  Section~\ref{sec:theory} gives the
candidate approximation and robust MaxWeight stability theorem.
Section~\ref{sec:implementation} describes the Scheduleurm implementation and the
measured feasible-family construction.  Section~\ref{sec:experiments} reports
the replay, live-sanity, slack-accounting, production-coverage, and oracle-trace
evidence.  Section~\ref{sec:related} positions the work relative to GPU cluster
schedulers, queueing-network control, and learning-while-scheduling.
Section~\ref{sec:discussion} states the limitations and the remaining calibration
requirements.

\section{Model}
\label{sec:model}

\subsection{Notation}
\label{subsec:notation}

Table~\ref{tab:notation} fixes the symbol ledger used throughout the paper.
Each mathematical symbol is reserved for the meaning shown in the table; in
particular, \(Q\) is the queue vector, \(q\) is an arbitrary nonnegative support
weight, \(K_t\) is switching cost, and \(G_t\) is the risk or interference
penalty.

\begin{table}[t]
\TABLE{Notation ledger.\label{tab:notation}}{
\small
\begin{tabular}{@{}p{0.27\textwidth}p{0.62\textwidth}@{}}
\toprule
Symbol & Meaning \\
\midrule
\(\mathcal I,i\) & finite set of job classes and a class index \\
\(t,Q_i(t),Q(t)\) & slot, class-\(i\) backlog, and queue vector \\
\(a,a_t\) & configuration action and chosen action \\
\(\Afull,\Acand\) & full action family and fixed candidate family \\
\(\Acand_t,\Ameas\) & statewise candidate family and measured action slice \\
\(Z_t,\omega_t,A_i(t),S_i(t)\) & regime, exogenous randomness, class-\(i\) arrival, and class-\(i\) service \\
\(\mu^z(a),\mu(a),\lowermu_t(a)\) & regime service mean, main-theorem service mean, and conservative lower-service estimate \\
\(\mathcal C(\mathcal A),\Lambda(\mathcal A),\Delta(\mathcal A),x_a\) & service convex hull, downward-closed load region, probability simplex, and stationary mixing mass on action \(a\) \\
\(H_{\mathcal A}(q),q\) & support function and nonnegative support vector \\
\(\Phi,d_{\Phi},L,\rho,\pi\) & finite action features, fabric metric, service Lipschitz envelope, cover radius, and candidate projection \\
\(K_t(a),G_t(a),P_0,\beta\) & switching cost, risk/interference penalty, additive penalty bound, and queue-scaled penalty coefficient \\
\(\delta,\epsilon_{\mathrm{cand}},\epsilon_{\mathrm{est}},\alpha_0,\alpha_1,\eta\) & full-region slack, candidate loss, estimation loss, additive oracle error, queue-scaled oracle error, and residual drift margin \\
\(B,N,V\) & second-moment bound, finite-set threshold, and quadratic Lyapunov function \\
\(k,n,M_k(n),F_k(n)\) & measured co-location profile, task-list size, all-job makespan proxy, and mean-flow proxy \\
\bottomrule
\end{tabular}
}{The notation avoids reusing one symbol for multiple meanings.  Acronyms are
expanded at first textual use and then used consistently.}
\end{table}

\subsection{Configuration actions}
\label{subsec:configuration-actions}

Let \(\mathcal I\) be a finite set of job classes.  Queue \(Q_i(t)\in\mathbb Z_+\)
counts the amount of unfinished work of class \(i\) at the beginning of slot
\(t\).  A global configuration action \(a\in\Afull\) specifies which jobs are
served, which resource configuration each job receives, and where each job is
placed.  In a GPU cluster this includes, for example, the number of jobs sharing
a GPU, CPU worker counts, memory pressure buckets, and host or device placement.
The model intentionally treats these as one action rather than as independent
per-job choices.

Given a regime \(z\) and exogenous randomness \(\omega\), action \(a\) generates
a nonnegative service vector \(S(a,z,\omega)\in\Real_+^{|\mathcal I|}\).  We write
\(\mu^z(a)=\E[S(a,z,\omega)]\).  For the main theorem we first suppress \(z\)
and use \(\mu(a)\); uniform-in-regime and average-regime extensions require
additional switching assumptions and are not used as main empirical claims.

The queue dynamics are
\begin{equation}
  Q_i(t+1)=\left[Q_i(t)-S_i(t)\right]^+ + A_i(t),\qquad i\in\mathcal I,
  \label{eq:queue-dynamics}
\end{equation}
where \(A_i(t)\) is the class-\(i\) arrival in slot \(t\).  In the finite-support
specialization used by the formal proof, the conditional distribution of
\((A(t),S(t))\) depends on \(Q(t)\) through a finite sample space.  In the more
general theorem it is enough to assume bounded conditional second moments.

\subsection{Capacity and support functions}
\label{subsec:capacity-support}

For a finite action family \(\mathcal A\), define the service region
\[
  \mathcal C(\mathcal A)
  =
  \left\{\sum_{a\in\mathcal A}x_a\mu(a):x\in\Delta(\mathcal A)\right\},
\]
and its downward closure
\[
  \Lambda(\mathcal A)
  =
  \left\{ \lambda\in\Real_+^{|\mathcal I|}:
    \exists v\in\mathcal C(\mathcal A)\ \text{with}\ \lambda\le v
  \right\}.
\]
The support function in the nonnegative orthant is
\[
  H_{\mathcal A}(q)=\max_{a\in\mathcal A}q^\top \mu(a),
  \qquad q\in\Real_+^{|\mathcal I|}.
\]
Here \(\Delta(\mathcal A)\) is the probability simplex over the action family
\(\mathcal A\), and \(x_a\) is the stationary mixing mass assigned to action
\(a\).  The vector \(q\) is a nonnegative support weight; it is not the queue
state unless explicitly specialized to \(Q(t)\).
If \(\lambda\) has coordinate slack \(\delta>0\) in the full region, then there
is a stationary mixture \(x\in\Delta(\Afull)\) such that
\(\lambda_i+\delta\le \sum_a x_a\mu_i(a)\) for every \(i\).  Consequently,
\begin{equation}
  q^\top\lambda+\delta\normone{q}\le H_{\Afull}(q),
  \qquad q\ge0.
  \label{eq:support-slack}
\end{equation}
This elementary support implication is the bridge from capacity to MaxWeight
drift; it is not an operational stability ``if and only if'' by itself.

\subsection{Candidate families and fabric metrics}
\label{subsec:candidate-cover}

The implemented scheduler optimizes over \(\Acand\subseteq\Afull\).  To avoid
turning this computational restriction into an unaccounted model change, define
a finite feature map \(\Phi:\Afull\to\Real^d\) and weighted \(\ell_1\) fabric metric
\[
  d_{\Phi}(a,a')=\sum_{r=1}^d w_r|\Phi_r(a)-\Phi_r(a')|.
\]
The features may include co-location profile, node bucket, CPU bucket, memory
pressure bucket, placement locality, and other scheduler-visible fabric
attributes.  We say that \(\Acand\) is a \(\rho\)-cover of \(\Afull\) if for each
full action \(a\) there is a candidate action \(\pi(a)\in\Acand\) with
\(d_{\Phi}(a,\pi(a))\le\rho\).  If
\(\|\mu(a)-\mu(a')\|_{\infty}\le L d_{\Phi}(a,a')\), then
\begin{equation}
  H_{\Afull}(q)\le H_{\Acand}(q)+L\rho\normone{q},
  \qquad q\ge0.
  \label{eq:cover-support-loss}
\end{equation}
This statement is used as a theorem assumption unless a perturbation-profiling
table certifies \(\Phi,L,\rho\).  For exact measured finite slices in which the
candidate family enumerates the measured actions, we set \(\rho=0\).

\section{Robust Candidate MaxWeight}
\label{sec:theory}

\subsection{Policy}
\label{subsec:policy}

Let \(\lowermu_t(a)\) be a conservative lower-service estimate for candidate
action \(a\) at time \(t\).  Let \(K_t(a)=K(a_{t-1},a)\) model switching or
rollback cost from the previous action to \(a\), and let \(G_t(a)\) model risk
or interference penalties.  The exact robust candidate MaxWeight rule is
\begin{equation}
  \begin{aligned}
  a_t\in\arg\max_{a\in\Acand_t}
  \bigl\{&
    Q(t)^\top\lowermu_t(a)\\
    &{}-K_t(a)-G_t(a)
  \bigr\}.
  \end{aligned}
  \label{eq:robust-policy}
\end{equation}
For practical schedulers, the selected action may be only approximately optimal.
We assume the queue-scaled oracle condition relative to
\eqref{eq:robust-policy},
\begin{equation}
  \begin{aligned}
  &\max_{a\in\Acand_t}
  \left\{Q^\top\lowermu_t(a)-K_t(a)-G_t(a)\right\}\\
  &\quad{}-
  \left\{Q^\top\lowermu_t(a_t)-K_t(a_t)-G_t(a_t)\right\}\\
  &\le \alpha_0+\alpha_1\normone{Q}.
  \end{aligned}
  \label{eq:oracle-error}
\end{equation}
Exact maximization is the special case \(\alpha_0=\alpha_1=0\).

\subsection{Main stability certificate}
\label{subsec:main-theorem}

\begin{theorem}[Robust candidate MaxWeight stability certificate]
\label{thm:main}
Consider a finite system with job classes \(\mathcal I\) and queue dynamics
\eqref{eq:queue-dynamics}.  Suppose:
\begin{enumerate}[leftmargin=18pt,itemsep=2pt]
\item \(\lambda\) has full-action support slack \(\delta\), i.e.,
      \eqref{eq:support-slack} holds.
\item The candidate family has support loss at most
      \(\epsilon_{\mathrm{cand}}\normone{q}\), where
      \(\epsilon_{\mathrm{cand}}=L\rho\) under the fabric-cover condition
      \eqref{eq:cover-support-loss}.
\item The lower-service map has support error at most
      \(\epsilon_{\mathrm{est}}\normone{q}\):
      \(H_{\Acand}(q)\le\max_{a\in\Acand}q^\top\lowermu(a)
      +\epsilon_{\mathrm{est}}\normone{q}\).
\item Penalties satisfy
      \(0\le K_t(a_t)+G_t(a_t)\le P_0+\beta\normone{Q(t)}\).
\item The selected action satisfies the approximate oracle condition
      \eqref{eq:oracle-error}.
\item Conditional moments satisfy
      \(\E[A_i(t)\mid Q(t)=Q]\le\lambda_i\),
      \(\lowermu_i(a_t)\le\E[S_i(t)\mid Q(t)=Q]\), and
      \[
        \E\!\left[\frac12\sum_i(A_i(t)^2+S_i(t)^2)\mid Q(t)=Q\right]\le B.
      \]
\end{enumerate}
If
\[
  \eta=\delta-(\epsilon_{\mathrm{cand}}+\epsilon_{\mathrm{est}}+\beta+\alpha_1)>0,
\]
then for \(V(Q)=\frac12\sum_i Q_i^2\) and
\(\Delta V(t)=V(Q(t+1))-V(Q(t))\),
\[
  \begin{aligned}
  \E[\Delta V(t)\mid Q(t)=Q]
  &\le B+P_0+\alpha_0\\
  &\quad{}-\eta\normone{Q}.
  \end{aligned}
\]
Consequently, for any \(N\) satisfying
\[
  B+P_0+\alpha_0+\gamma \le \eta N
\]
for some \(\gamma>0\), the process has a negative expected drift outside
\(\{Q:\normone{Q}\le N\}\), and hence satisfies a finite-set Foster recurrence
certificate.
\end{theorem}

\noindent
The Lean 4 proof artifact \citep{deMoura2021lean} proves the fixed-family
theorem and a statewise version in which the feasible family \(\Acand_t\)
depends on the queue state, available jobs, and measured capacity-boundary rows.
It also proves a bounded finite-support specialization in which \(B\) is
constructed from coordinate sample bounds.  In the artifact map, the
paper-facing results are indexed as the calibrated approximate-oracle theorem
and the statewise calibrated approximate-oracle theorem.

\subsection{Proof decomposition}
\label{subsec:proof-decomposition}

The proof has four parts.  First, full-action slack implies the support bound
\eqref{eq:support-slack}.  Second, the fabric cover transfers support from
\(\Afull\) to \(\Acand\) with loss \(L\rho\normone{q}\).  Third,
lower-service error, penalties, and approximate optimization are subtracted from
the same slack budget.  Fourth, the standard quadratic Lyapunov inequality
\[
  \begin{aligned}
  &V([Q-S]^+ + A)-V(Q)\\
  &\quad\le \frac12\sum_i(A_i^2+S_i^2)\\
  &\qquad{}+Q^\top A-Q^\top S
  \end{aligned}
\]
converts the deterministic support margin into conditional negative drift.  The
formalization keeps the operational necessity direction separate: an
``operationally stabilizes'' predicate is load-certified by a model that
explicitly encodes the mean arrival vector, and zero-slack necessity follows
from a conservation-law argument rather than from a definition.

\section{Scheduleurm Instantiation}
\label{sec:implementation}

\subsection{Measured action slices}
\label{subsec:measured-slices}

Scheduleurm is an existing scheduler for local and remote research jobs.  The
new algorithmic layer is separated from the live scheduler: service profiling,
replay, SOTA-style baselines, slack accounting, and oracle audits live under
\texttt{simulation/} and \texttt{algorithm/experiments/}.  The live scheduler can
select between legacy and algorithmic modes, but the theorem-facing evidence is
constructed from explicit artifacts rather than from undocumented changes to
the production dispatch loop.

Table~\ref{tab:bucket-slices} summarizes the current measured finite slices.
A capacity boundary is a measured or live sanity profile that is not usable as a
robust feasible action.  Once profile \(k\) is a boundary, profile \(k\) and all
higher profiles are excluded from the theorem-facing candidate family for that
workload bucket.

\begin{table}[t]
\TABLE{Current measured finite slices.\label{tab:bucket-slices}}{
\begin{tabular}{lrrrr}
\toprule
Bucket & Feasible profiles & Boundary & Candidate & Legacy \\
\midrule
q00 \texttt{light\_control\_local} & 1--13 & 14 & 13 & 1 \\
q01 \texttt{gpu\_heavy\_jax\_matmul} & 1--8 & none & 4 & 3 \\
q10 \texttt{cpu\_heavy\_local\_bench} & 1--9 & 10 & 8 & 9 \\
q11 \texttt{hybrid\_rl\_resac\_ant} & 1--9 & 10 & 2 / 3 & 5 \\
\bottomrule
\end{tabular}
}{q11 uses profile 2 in the standalone task list and profile 3 in the mixed portfolio.
The q11 profile-10 replay point from earlier measurements is historical only:
fresh live sanity showed runtime out-of-memory (OOM) at profile 10, making it
the first robust capacity boundary for the current node bucket.}
\end{table}

The replay policy uses a support-first selector.  For a single workload, it
computes a deterministic all-job makespan proxy \(M_k(n)\) for each measured
profile \(k\), where \(n\) is the task-list size, keeps profiles within a
makespan guard, and then minimizes a deterministic mean-flow proxy \(F_k(n)\)
within that guard.  For the mixed portfolio, it first preserves the global
makespan/support objective and then spends the remaining guard on weighted mean
flow.  This implementation matches the approximate-oracle theorem: the support
guard corresponds to the
\(\alpha_0+\alpha_1\normone{Q}\) term, and the finite-batch delay tie-break is a
bounded or queue-scaled penalty.

\subsection{Auditable theorem objects}
\label{subsec:auditable-objects}

The Scheduleurm artifact distinguishes three populations.  The explicit
task-list replay population is used for performance comparisons.  The
completed-active production population is used for a service-map theorem oracle
bridge.  The raw scheduler history is treated only as telemetry: it contains
cancelled, attempted-only, benchmark, and externally adopted jobs and is not a
clean arrival stream for the theorem.

The current completed-active Scheduleurm-controlled production population has
3437 mapped records out of 3437, positive mapped-slice capacity slack, and a
service-map robust MaxWeight lower-service oracle bridge with
\(\alpha_0=\alpha_1=0\).  A separate live trace audit closes one emitted
candidate-family trace on two local CPU control-plane slots.  This proves the
trace/enrichment/audit pipeline is executable; it does not certify every future
dispatch without rerunning the audit.

\section{Computational Evidence}
\label{sec:experiments}

\subsection{Explicit task-list replay}
\label{subsec:task-list-replay}

The main benchmark uses one explicit task list per quadrant and one mixed
portfolio.  All policies receive the same jobs, total work units, resource-pool
counts, arrival times, and measured service cache.  Static arrivals are used for
the first comparison because all-job completion time is then unambiguous.
Poisson and load-sweep traces are implemented but are not used as the primary
claim in this draft.

Table~\ref{tab:legacy-results} compares the current candidate policy with the
legacy Scheduleurm caps.  The metric is completion of all jobs in the task list,
not the time to finish a single job.

\begin{table}[t]
\TABLE{Candidate policy versus legacy Scheduleurm on explicit task lists.\label{tab:legacy-results}}{
\begin{tabular}{l l r r r}
\toprule
Task list & Candidate profile & Makespan & Mean flow & p90 flow \\
\midrule
q00 light control & \texttt{light=13} & \(11.825\times\) & \(11.758\times\) & \(11.814\times\) \\
q01 GPU-bound compute & \texttt{gpu=4} & \(1.079\times\) & \(1.045\times\) & \(1.085\times\) \\
q10 CPU-host bound & \texttt{cpu=8} & \(1.510\times\) & \(1.534\times\) & \(1.534\times\) \\
q11 CPU/GPU coupled & \texttt{hybrid=2} & \(1.154\times\) & \(1.176\times\) & \(1.156\times\) \\
Mixed portfolio & \texttt{cpu=8,gpu=1,hybrid=3} & \(1.161\times\) & \(1.269\times\) & \(1.186\times\) \\
\bottomrule
\end{tabular}
}{Ratios are legacy divided by candidate, so larger is better for Scheduleurm candidate.
Each row uses the current robust feasible family; q11 profile 10 and above are excluded.
p90 is the 90th percentile of job flow time.}
\end{table}

\subsection{SOTA-style policy baselines}
\label{subsec:sota-style}

No single external scheduler covers all four Scheduleurm quadrants and the
mixed portfolio without changing the workload model.  We therefore compare to
SOTA-style policy semantics on the same measured service cache: a
throughput-table goodput policy representing Gavel, Pollux, and Sia-like
schedulers; a finite-batch delay oracle representing shortest remaining
processing time (SRPT), Gittins, and shortest expected remaining processing time
(SERPT)-style policies \citep{scully2020simple}; and an interference guard
representing IADeep/Salus-like co-location control.  This is a policy-semantics
comparison, not a direct binary execution of those systems.

Table~\ref{tab:sota-matrix} reports the fixed-policy matrix.  The candidate is
not Pareto-dominated on total makespan and job-weighted mean flow.  The matrix
also shows the nature of the tradeoff: throughput-table policies can be
slightly faster on aggregate makespan, while delay-oriented policies can be
slightly better on mean flow.

\begin{table}[t]
\TABLE{Fixed SOTA-style policy matrix on identical measured service cache.\label{tab:sota-matrix}}{
\begin{tabular}{l r r r r}
\toprule
Policy & Jobs & Sum makespan & Mean flow & Candidate vs. policy \\
\midrule
Candidate & 1376 & 81818.460 & 6060.620 & \(1.0000,1.0000\) \\
Throughput table & 1376 & 81793.703 & 6067.049 & \(0.9997,1.0011\) \\
Delay oracle & 1376 & 81988.550 & 6052.181 & \(1.0021,0.9986\) \\
Interference guard & 1376 & 81964.959 & 6055.165 & \(1.0018,0.9991\) \\
Quadrant composite & 1376 & 81964.959 & 6055.165 & \(1.0018,0.9991\) \\
\bottomrule
\end{tabular}
}{The last column gives candidate-versus-policy ratios for (sum makespan, job-weighted
mean flow).  Values above one favor the candidate.  No row dominates the
candidate on both objectives.}
\end{table}

On the q01 GPU-heavy list, the throughput endpoint uses profile 8 and improves
all-job time by about 1.6\% relative to the guarded candidate, but the candidate
improves mean flow by about 11.6\% relative to that endpoint.  The delay
endpoint uses profile 1 and improves mean flow, but loses all-job time.  This is
the intended Pareto geometry rather than a hidden failure of the module.

\subsection{Theorem-condition and live-sanity checks}
\label{subsec:condition-checks}

The measured portfolio finite slice also instantiates the slack inequality in
Theorem~\ref{thm:main}.  With load
\(\lambda_i=0.8\,\mu_i(a_{\mathrm{selected}})\), the candidate family equals
the measured action slice, so \(L\rho=0\).  The resulting certificate is
\[
  \begin{aligned}
  \delta &= 0.066921842,\\
  \epsilon_{\mathrm{est}}=\beta=\alpha_1 &= 0,\\
  \eta &= 0.066921842.
  \end{aligned}
\]
The finite-support second-moment bound is \(B=6361.354514\), and the resulting
finite-set threshold is \(N=95072\).  This is a theorem-condition certificate for
a declared finite-support load model; it is not a claim that the static replay
trace is itself a stationary arrival process.

The replay model was checked against fresh progress-window live measurements.
For q01 profile 4 on two GPUs, replay predicted a makespan of 1661.313 seconds,
whereas the fresh live proxy was 1652.248 seconds, a relative error of
0.005457.  For q11 profile 2, replay predicted 38382.365 seconds and the fresh
live proxy was 37760.000 seconds, a relative error of 0.016215.  For the mixed
portfolio, composed live slices gave relative errors of 0.007349 on makespan,
0.005417 on mean flow, and 0.007964 on p90 flow.  These are progress-window
sanity checks; the long jobs were not all run to natural completion.

\subsection{Formal and production artifacts}
\label{subsec:formal-production}

The Lean proof artifact is submitted separately as a consolidated file
\texttt{ScheduleurmUpload.lean}.  The checked artifact has Secure Hash
Algorithm 256-bit (SHA-256) hash
\begin{center}
{\scriptsize\ttfamily af79be4416e4c4add0fe41663fc0927af7058fe04412908b6688a8409227f01b}.
\end{center}
The commands \texttt{lake build Scheduleurm} and
\texttt{lake env lean ScheduleurmUpload.lean} both pass, and a search for
\texttt{sorry}, \texttt{admit}, and \texttt{axiom} returns no matches.  The
theorem names used in this paper are searchable in the upload file.

For production data, the current theorem-facing population is the
completed-active Scheduleurm-controlled set.  It has 3437 strict mapped records
out of 3437, with mapped-slice slack \(9.123\times10^{-6}\).  The service-map
oracle bridge passes with \(\alpha_0=\alpha_1=0\).  The raw 30-day history has
unmapped records and is not used as the theorem population.  This distinction is
important: the artifact proves coverage of a controlled completed-active
population, not closure of every raw or future scheduler row.

\section{Related Work}
\label{sec:related}

\subsection{Deep-learning and heterogeneous cluster schedulers}
\label{subsec:related-systems}

Goodput-oriented deep-learning schedulers such as Gavel, Pollux, and Sia use
measured or modeled throughput tables to allocate heterogeneous accelerators
\citep{narayanan2020gavel,qiao2021pollux,sia2023}.  Fairness-oriented systems
such as Themis study efficient sharing under fairness constraints
\citep{mahajan2020themis}.  GPU sharing systems such as Gandiva, Salus, and
IADeep expose or exploit fine-grained co-location and interference behavior
\citep{xiao2018gandiva,yu2020salus,iadeep2023}.  These systems motivate our
measured-service replay baselines, but our theoretical question is different:
we ask how a scheduler can restrict a combinatorial configuration action space
to a candidate family while preserving a queueing-stability certificate with
explicit loss terms.

Hybrid CPU/GPU placement and cluster allocation systems such as AlloX and
related heterogeneous placement work also emphasize that accelerator scheduling
cannot be reduced to GPU count alone \citep{le2020allox,carvalho2024workload}.
Scheduleurm follows the same empirical lesson.  Its q00, q10, and q11 buckets
are declared local or node-bucket service certificates, not claims that all
CPU-heavy or RL workloads have the same service curve.

\subsection{Queueing-network control and learning}
\label{subsec:related-queueing}

MaxWeight and backpressure policies are classical tools for stabilizing
constrained queueing systems under capacity-region slack
\citep{tassiulas1992stability,stolyar2004maxweight,busic2013stability}.  Our
proof follows this lineage but adds three implementation-facing terms:
candidate-set approximation, lower-service estimation, and approximate-oracle
loss.  Work on state-dependent processor sharing and workload-dependent
admission control studies service rates that depend on the number or composition
of active jobs \citep{gupta2022state,bekker2006admission,altman2006dps}.  The
Scheduleurm co-location curves are a systems instance of the same phenomenon.

Learning-while-scheduling work, including MaxWeight with discounted upper
confidence bound (UCB) and learning-based queueing controls, studies unknown
service statistics under online exploration
\citep{yang2023maxweight,liu2020rlqn,dai2022deep}.  We do
not claim a complete online learning theorem for the current Scheduleurm
sampler.  The formal artifact proves conditional high-probability lifting for
active-bucket certificates, but the main empirical claims use measured
service-cache certificates.

\section{Discussion and Limitations}
\label{sec:discussion}

The main contribution of this work is a bridge between Operations Research
(OR)-style queueing certificates and a concrete scheduler artifact.  The theory
does not shrink the full action space into the implemented candidate set.
Instead, it states the loss incurred by that restriction and pays it from
capacity slack.  The Scheduleurm implementation then makes the proof objects
auditable: measured lower-service rows, capacity boundaries, explicit task
lists, finite-slice slack accounting, and lower-service oracle traces are all
materialized as artifacts.

Several limitations remain.  First, the strongest empirical claim is an exact
measured finite-slice claim.  The Lean theorem supports generalized fabric-cover
calibration, but a broad claim requires a perturbation-profiling table with
features \(\Phi\), weights, projection \(\pi\), cover radius \(\rho\), service
sensitivity envelope \(L\), and failure probability if the envelope is
statistical.  Second, the SOTA comparisons are replayed policy semantics on the
same measured Scheduleurm service cache.  They are not direct executions of
Gavel, Pollux, Sia, IADeep, or Salus.  Third, the live oracle trace currently
certifies one emitted candidate-family trace; future online theorem-grade
claims require the trace/enrichment/audit pipeline to be rerun.  Fourth, q00,
q10, and q11 are declared local or node-bucket certificates.  They should not be
generalized to all CPU-heavy, data-loader-heavy, RL, or BAPR deployments without
replication on those buckets.

These limitations are not incidental.  They are the boundary that makes the
claims falsifiable.  The next empirical step is to replace the exact finite-slice
claim by a calibrated fabric-cover table and to run longer subcritical
arrival-process experiments.  The next theoretical step is to connect the
conditional active-bucket learning and hidden-regime extension theorems to the
actual Scheduleurm sampling and detection mechanisms.

\section{Conclusion}
\label{sec:conclusion}

We formulated heterogeneous compute scheduling as a configuration-action
stochastic processing network and proved a robust candidate MaxWeight stability
certificate with explicit slack accounting.  The proof identifies precisely how
candidate approximation, lower-service estimation, penalties, and approximate
optimization consume capacity slack.  In Scheduleurm, measured finite service
slices, explicit task-list replay, live progress-window checks, production
coverage, and oracle-trace audits instantiate the theorem objects under a
bounded claim scope.  The current evidence supports a strong but limited
conclusion: on declared measured slices, Scheduleurm's robust candidate policy
improves over its legacy caps and is not Pareto-dominated by SOTA-style replay
baselines, while the formal artifact verifies the stability spine without
unproved Lean placeholders.

\bibliographystyle{informs2014}
\bibliography{references}

\end{document}
